\documentclass{article} % This command is used to set the type of
% document you are working on such as an article, book, or presenation

\usepackage{geometry} % This package allows the editing of the page layout
\usepackage{amsmath}  % This package allows the use of a large range
% of mathematical formula, commands, and symbols
\usepackage{graphicx}  % This package allows the importing of images
\usepackage{amssymb}
\usepackage{amsfonts}
\newcommand{\question}[2][]{
  \begin{flushleft}
    \textbf{Question #1}: \textit{#2}

\end{flushleft}}
\newcommand{\sol}{\textbf{Solution}:} %Use if you want a boldface solution line
\newcommand{\maketitletwo}[2][]{
  \begin{center}
    \Large{\textbf{Assignment #1}

    Topology} % Name of course here
    \vspace{5pt}

    \normalsize{N. Ohayon Rozanes  % Your name here

    \today}        % Change to due date if preferred
    \vspace{15pt}

\end{center}}
\begin{document}
\maketitletwo[7]  % Optional argument is assignment number
%Keep a blank space between maketitletwo and \question[1]

\question[1]{}
\begin{enumerate}
  \item
    A linear function that maps between $(a, b)$ to $(-1,1)$ would
    suffice, this can be achieved with the mapping

    \[
      f(x) = \frac{2}{b-a}x - \frac{a+b}{b-a}
    \]

    First we verify that indeed $f((a,b)) = (-1,1)$,
    \begin{align*}
      a < x < b &\iff \frac{2a}{b-a} < \frac{2}{b-a}x < \frac{2b}{b-a} \\
      &\iff \frac{2a - b - a}{b-a} < \frac{2}{b-a}x -
      \frac{b+a}{b-a}< \frac{2b
      - b - a}{b-a} \\
      &\iff \frac{a - b}{b-a} < \frac{2}{b-a}x - \frac{b+a}{b-a}<
      \frac{b-a}{b-a} \\
      &\iff -1 < \frac{2}{b-a}x - \frac{b+a}{b-a}<
      1 \\
      &\iff -1 < f(x)<
      1 \\
    \end{align*}

    We verify that the function is bijective, continuous and with a
    continuous inverse:

    A function is bijective if $\forall y \in Y, \exists!x \in X, y =
    f(x)$, take any $y \in Y$, then $y = f(x) = ax + b \implies
    \frac{y-b}{a} = x$, which in this case
    $$
    x = \frac{b-a}{2}y + \frac{a+b}{2}
    $$

    As demonstrated for any $y \in  (-1,1) \implies x\in(a,b)$, and this $x$ is
    unique, therefore the function is bijective from $(a,b) \to (-1,1)$.

    To demonstrate that this function and it's inverse are continuous,
    consider that for any open set $A \subset (-1,1)$, since open
    intervals form a basis in $\mathbb{R}$,

    \[
      A = \bigcup (a_i, b_i) \implies f^{-1}(A) = \bigcup f^{-1}(a_i, b_i)
    \]

    It suffices to show that for any interval $(a_i, b_i) \subset (-1,1)$,
    $f^{-1}(a_i, b_i)$ is open.

    \begin{align*}
      a_i < f(x) < b_i &\iff a_i < \frac{2}{b-a}x - \frac{b+a}{b-a} < b_i \\
      &\iff \frac{b-a}{2}a_i + \frac{b+a}{2} < x < \frac{b-a}{2}b_i +
      \frac{b+a}{2} \\
      &\iff f^{-1}(a_i) < x < f^{-1}(b_i)
    \end{align*}
    Which is the open interval $(f^{-1}(a_i), f^{-1}(b_i)) \subset
    (a,b)$, since the function is bijective, and the above shows
    equivalency, the same argument holds for $f^-1$. That is to say, that
    both $f$ and $f^{-1}$ are continuous. As such, $f$ is bijective,
    continuous, and with a continuous inverse, as such, it is a homeomorphism.
  \item

    % this might be too lazy, largely the proof is that tan on the
    % domain continuous since it maps intervals to intervals,this can
    % be shown from inequality equivalences
    % it is bijective as you take limits

    % then arctan follows since tan is increasing and continuous
    Using the previous part, and the fact that the composition of
    bijective, continuous mappings is bijective and continuous, it
    suffices to find a homeomorphism $g: (-1, 1) \to (-\infty, \infty)$
    is a homeomorphism, since then $g
    \circ f$ would be a homeomorphism from $(a, b) \to (-\infty,
    \infty)$. For this $g(x) = \tan(\frac{\pi}{2}x)$ satisfies the
    requirements, since it is bijective, is continuous on the domain,
    and has a continuous inverse $\frac{\pi}{2}\arctan(x)$. As such,
    the function
    $$
    \tan\left(\frac{\pi}{2}\left(\frac{b-a}{2}x + \frac{b+a}{2} \right)\right)
    $$
    is a homeomorphism
\end{enumerate}
\question[2]{}

Using the stereographic projection, which draws a line from the
removed point $(1,0)$ to every other point on the circle $(x,y)$, and
sets $f(x,y)$ to where it crosses the $y$ axis, we obtain the
following explicit formula:

$$
f(x,y) = \frac{y}{1-x}
$$

We now seek to verify that it is bijective, that it is continuous,
and that it's inverse is continuous.

To show bijectivity, take any $b \in \mathbb{R}$ then we can solve
the following equations for $x,y \in C\setminus\{(0,1)\}$

\[
  b = \frac{y}{1-x}
\]

\[
  x^2 + y^2 = 1
\]

It follow from the first that $y = (1-x)b$, substituting that into
the second yields
\[
  x^2 + ((1-x)b)^2 = 1 \implies x^2(1+b^2) - 2b^2x + (b^2-1) = 0
\]

As a quadratic equation, this has two solutions, however, note that
$x = 1$ satisfies the above equation
\[
  1+b^2 - 2b^2 + b^2-1 = 0
\]

Therefore, there is only one $x \in C\setminus\{(0,1)\}$ which
satisifies the conditions, and it follows that there exists only one
$y$ which does as well. As such, for any $b\in\mathbb{R}$ there
exists only one $(x,y) \in C\setminus\{(0,1)\}$ such that $f(x,y) =
b$, and therefore the function is bijective. From this we can also
compute the explicit inverse function, using the quadratic formula

\[
  x = \frac{2b^2 - \sqrt{4b^4 - 4(1+b^2)(b^2-1)}}{2(1+b^2)} =
  \frac{2b^2 - 2}{2(1+b^2)} = \frac{b^2 -1}{b^2+1}
\]
Therefore

$$
f^{-1}(b) = \left(\frac{b^2-1}{b^2+1}, \frac{2b}{b^2+1}\right)
$$

%% continuity proofs of various kinds proceed, I think the proof for
% 1/b^2 + 1 and b^2-1 are simple enough,
%% then composition of continuous + continuous components is
% continuous. For f I'm not sure how to do that, it probably follows
% pretty quickly with the subspace topology

\question[3]{}
To show that $\mathcal{I}$ is not connected, we find two disjoint
open sets $U, V$ in the subspace topology of $\mathcal{I}$ such that

$$
\mathcal{I} = U \cup V
$$

Take $U = (-\infty, 0) \cap \mathcal{I}$ and $V = (0,\infty) \cap
\mathcal{I}$, clearly, since $(-\infty, 0)$ and $(0, \infty)$ are in
the topology of $\mathbb{R}$, $U, V  \in \mathcal{T}_\mathcal{I}$, further

$$
U \cup V = (-\infty, \infty)  \cap \mathcal{I} = \mathcal{I} \quad U
\cap V = \emptyset
$$
As such, $U,V$ is a seperation of $\mathcal{I}$, and therefore
$\mathcal{I}$ is not connected.
\question[4]{}

Let $x \in [a,b]$, assume for the sake of contradiction that $x
\notin A$, then take $U = (-\infty, x)$ and $V = (x, \infty)$. $U,V$
are open in $\mathbb{R}$, so $U \cap A$ and $V \cap A$ are open in
the subset topology of $A$. Further more, since $x\in [a,b] \implies a
\leq x \leq b$, and $a,b \in A$, we have that at least $a \in U \cap
A$ and $b \in V \cap A$, which is to say neither sets are empty. Further

$$
(U \cap A) \cup (V \cap A)  = (U \cup V) \cap A = \mathbb{R}\setminus
\{x\} \cap A = A
$$

And also
$$
(U \cap A) \cap (V \cap A) = (U \cap V) \cap A  = \emptyset \cap A = \emptyset
$$

Therefore, $U\cap A$ and $V \cap A$ seperate $A$, however, $A$ is
connected, which is a contradiction. As such $\forall x \in [a,b],
x\in A \implies [a,b] \subset A$
\end{document}
